\subsection{\label{sec.cancel.moresub}More subtractive methods}

\subsubsection{\label{subsec.cancel.moresub.all-even}Sums with varying signs}

The ideas of subtraction and cancellation have several other applications in
combinatorics. One way to achieve a lot of cancellation is by summing the same
sums but with varying signs. For instance, when $P\left(  x\right)  $ is a
univariate polynomial, the sum $P\left(  x\right)  +P\left(  -x\right)  $
contains no odd powers of $x$, since all these odd powers cancel out (for
instance, $\left(  1+x\right)  ^{3}+\left(  1-x\right)  ^{3}=6x^{2}+2$).
Meanwhile, all even powers of $x$ that appear in $P\left(  x\right)  $ will
appear in $P\left(  x\right)  +P\left(  -x\right)  $ with their coefficients
doubled. Likewise, for a polynomial $P\left(  x,y\right)  $ in two variables,
we can \textquotedblleft disable\textquotedblright\ all odd powers of $x$ by
taking $P\left(  x,y\right)  +P\left(  -x,y\right)  $; moreover, by taking the
sum $P\left(  x,y\right)  +P\left(  -x,y\right)  +P\left(  x,-y\right)
+P\left(  -x,-y\right)  $, we can \textquotedblleft disable\textquotedblright%
\ all monomials $x^{i}y^{j}$ for which at least one of $i$ and $j$ is odd,
leaving only the monomials $x^{i}y^{j}$ with $i\equiv j\equiv
0\operatorname{mod}2$ intact.

An impressive example of how this trick can be used is the proof of the
following enumerative result:

\begin{theorem}
\label{thm.cancel.all-even}Let $n\in\mathbb{N}$ and $d\in\mathbb{N}$. An
$n$-tuple $\left(  x_{1},x_{2},\ldots,x_{n}\right)  \in\left[  d\right]  ^{n}$
is said to be \emph{all-even} if each element of $\left[  d\right]  $ occurs
an even number of times in this $n$-tuple (i.e., if for each $k\in\left[
d\right]  $, the number of all $i\in\left[  n\right]  $ satisfying $x_{i}=k$
is even). For example, the $4$-tuple $\left(  1,4,4,1\right)  $ and the
$6$-tuples $\left(  1,3,3,5,1,5\right)  $ and $\left(  2,4,2,4,3,3\right)  $
are all-even, while the $4$-tuples $\left(  1,2,2,4\right)  $ and $\left(
2,4,6,4\right)  $ are not.

Then, the number of all all-even $n$-tuples $\left(  x_{1},x_{2},\ldots
,x_{n}\right)  \in\left[  d\right]  ^{n}$ is
\[
\dfrac{1}{2^{d}}\sum_{k=0}^{d}\dbinom{d}{k}\left(  d-2k\right)  ^{n}.
\]

\end{theorem}

This formula entails that the sum $\sum_{k=0}^{d}\dbinom{d}{k}\left(
d-2k\right)  ^{n}$ is nonnegative and divisible by $2^{d}$. This would not be
obvious without the theorem! The presence of addends with different signs (at
least for $n$ odd) appears to make a direct combinatorial proof impossible. We
will instead use the cancellation trick we mentioned above. (Our proof is
taken from \cite[solution to Exercise 7]{18f-hw4s}.)

As an introductory example, let us take a look at the case $n=4$ and $d=2$.
Then, Theorem \ref{thm.cancel.all-even} is saying that the number of all
all-even $4$-tuples $\left(  x_{1},x_{2},x_{3},x_{4}\right)  \in\left[
2\right]  ^{4}$ is
\begin{align*}
\dfrac{1}{2^{2}}\sum_{k=0}^{2}\dbinom{2}{k}\left(  2-2k\right)  ^{4}  &
=\dfrac{1}{4}\left(  \dbinom{2}{0}\cdot2^{4}+\dbinom{2}{1}\cdot0^{4}%
+\dbinom{2}{2}\cdot\left(  -2\right)  ^{4}\right) \\
&  =\dfrac{1}{4}\left(  2^{4}+2\cdot0^{4}+\left(  -2\right)  ^{4}\right)  .
\end{align*}
We claim that this can be seen as follows: For any integers $e_{1},e_{2}$, the
distributive law yields%
\begin{align*}
\left(  e_{1}+e_{2}\right)  ^{4}  &  =e_{1}e_{1}e_{1}e_{1}+e_{1}e_{1}%
e_{1}e_{2}+e_{1}e_{1}e_{2}e_{1}+\cdots\ \ \ \ \ \ \ \ \ \ \left(  \text{a
total of }16\text{ addends}\right) \\
&  =\sum_{\left(  x_{1},x_{2},x_{3},x_{4}\right)  \in\left[  2\right]  ^{4}%
}e_{x_{1}}e_{x_{2}}e_{x_{3}}e_{x_{4}}.
\end{align*}
Summing this over all four pairs $\left(  e_{1},e_{2}\right)  \in\left\{
1,-1\right\}  ^{2}$, we find%
\begin{align}
&  \sum_{\left(  e_{1},e_{2}\right)  \in\left\{  1,-1\right\}  ^{2}}\left(
e_{1}+e_{2}\right)  ^{4}\nonumber\\
&  =\sum_{\left(  e_{1},e_{2}\right)  \in\left\{  1,-1\right\}  ^{2}}%
\ \ \sum_{\left(  x_{1},x_{2},x_{3},x_{4}\right)  \in\left[  2\right]  ^{4}%
}e_{x_{1}}e_{x_{2}}e_{x_{3}}e_{x_{4}}\nonumber\\
&  =\sum_{\left(  x_{1},x_{2},x_{3},x_{4}\right)  \in\left[  2\right]  ^{4}%
}\ \ \sum_{\left(  e_{1},e_{2}\right)  \in\left\{  1,-1\right\}  ^{2}}%
e_{x_{1}}e_{x_{2}}e_{x_{3}}e_{x_{4}} \label{eq.thm.cancel.all-even.ex1}%
\end{align}
(here, we interchanged the summation signs). But now we can simplify the inner
sum $\sum_{\left(  e_{1},e_{2}\right)  \in\left\{  1,-1\right\}  ^{2}}%
e_{x_{1}}e_{x_{2}}e_{x_{3}}e_{x_{4}}$ rather conveniently:

\begin{itemize}
\item If the $4$-tuple $\left(  x_{1},x_{2},x_{3},x_{4}\right)  $ is all-even,
then all addends of this sum are $1$ (for example, if $\left(  x_{1}%
,x_{2},x_{3},x_{4}\right)  =\left(  1,2,2,1\right)  $, then $e_{x_{1}}%
e_{x_{2}}e_{x_{3}}e_{x_{4}}=e_{1}e_{2}e_{2}e_{1}=e_{1}^{2}e_{2}^{2}=1$ for all
$\left(  e_{1},e_{2}\right)  \in\left\{  1,-1\right\}  ^{2}$), and thus the
sum is equal to $4$.

\item If the $4$-tuple $\left(  x_{1},x_{2},x_{3},x_{4}\right)  $ is not
all-even, then the addends of this sum cancel out entirely (for example, if
$\left(  x_{1},x_{2},x_{3},x_{4}\right)  =\left(  1,2,1,1\right)  $, then
$e_{x_{1}}e_{x_{2}}e_{x_{3}}e_{x_{4}}=e_{1}e_{2}e_{1}e_{1}=e_{1}^{3}e_{2}$
changes sign whenever we flip the sign of $e_{1}$, so that the addends with
$e_{1}=1$ cancel out the addends with $e_{1}=-1$), and thus the sum is just
$0$.
\end{itemize}

Altogether, we thus see that each $\left(  x_{1},x_{2},x_{3},x_{4}\right)
\in\left[  2\right]  ^{4}$ satisfies
\[
\sum_{\left(  e_{1},e_{2}\right)  \in\left\{  1,-1\right\}  ^{2}}e_{x_{1}%
}e_{x_{2}}e_{x_{3}}e_{x_{4}}=%
\begin{cases}
4, & \text{if }\left(  x_{1},x_{2},x_{3},x_{4}\right)  \text{ is all-even};\\
0, & \text{if not.}%
\end{cases}
\]
Hence, (\ref{eq.thm.cancel.all-even.ex1}) simplifies to%
\begin{align*}
\sum_{\left(  e_{1},e_{2}\right)  \in\left\{  1,-1\right\}  ^{2}}\left(
e_{1}+e_{2}\right)  ^{4}  &  =\sum_{\left(  x_{1},x_{2},x_{3},x_{4}\right)
\in\left[  2\right]  ^{4}}%
\begin{cases}
4, & \text{if }\left(  x_{1},x_{2},x_{3},x_{4}\right)  \text{ is all-even};\\
0, & \text{if not}%
\end{cases}
\\
&  =4\cdot\left(  \text{\# of all all-even }4\text{-tuples }\left(
x_{1},x_{2},x_{3},x_{4}\right)  \in\left[  2\right]  ^{4}\right)  .
\end{align*}
We thus conclude that%
\begin{align*}
\left(  \text{\# of all all-even }4\text{-tuples }\left(  x_{1},x_{2}%
,x_{3},x_{4}\right)  \in\left[  2\right]  ^{4}\right)   &  =\dfrac{1}{4}%
\sum_{\left(  e_{1},e_{2}\right)  \in\left\{  1,-1\right\}  ^{2}}\left(
e_{1}+e_{2}\right)  ^{4}\\
&  =\dfrac{1}{4}\left(  2^{4}+2\cdot0^{4}+\left(  -2\right)  ^{4}\right)
\end{align*}
(since the sum $\sum_{\left(  e_{1},e_{2}\right)  \in\left\{  1,-1\right\}
^{2}}\left(  e_{1}+e_{2}\right)  ^{4}$ contains one addend with $e_{1}%
+e_{2}=2$, two addends with $e_{1}+e_{2}=0$, and one addend with $e_{1}%
+e_{2}=-2$). This is precisely the answer that Theorem
\ref{thm.cancel.all-even} gives us.

We shall now prove Theorem \ref{thm.cancel.all-even} by formalizing and
generalizing this argument. We begin with the obvious generalization of
(\ref{eq.thm.cancel.all-even.ex1}):

\begin{lemma}
\label{lem.cancel.all-even.l1}Let $n,d\in\mathbb{N}$. Then,%
\begin{align*}
&  \sum_{\left(  e_{1},e_{2},\ldots,e_{d}\right)  \in\left\{  1,-1\right\}
^{d}}\left(  e_{1}+e_{2}+\cdots+e_{d}\right)  ^{n}\\
&  =\sum_{\left(  x_{1},x_{2},\ldots,x_{n}\right)  \in\left[  d\right]  ^{n}%
}\ \ \sum_{\left(  e_{1},e_{2},\ldots,e_{d}\right)  \in\left\{  1,-1\right\}
^{d}}e_{x_{1}}e_{x_{2}}\cdots e_{x_{n}}.
\end{align*}

\end{lemma}

\begin{proof}
For each $\left(  e_{1},e_{2},\ldots,e_{d}\right)  \in\left\{  1,-1\right\}
^{d}$, we have $e_{1}+e_{2}+\cdots+e_{d}=\sum_{k=1}^{d}e_{k}$ and thus%
\begin{align*}
&  \left(  e_{1}+e_{2}+\cdots+e_{d}\right)  ^{n}\\
&  =\left(  \sum_{k=1}^{d}e_{k}\right)  ^{n}=\prod_{i=1}^{n}\ \ \sum_{k=1}%
^{d}e_{k}=\underbrace{\sum_{\left(  k_{1},k_{2},\ldots,k_{n}\right)
\in\left[  d\right]  \times\left[  d\right]  \times\cdots\times\left[
d\right]  }}_{=\sum_{\left(  k_{1},k_{2},\ldots,k_{n}\right)  \in\left[
d\right]  ^{n}}}\ \ \underbrace{\prod_{i=1}^{n}e_{k_{i}}}_{=e_{k_{1}}e_{k_{2}%
}\cdots e_{k_{n}}}\\
&  \ \ \ \ \ \ \ \ \ \ \ \ \ \ \ \ \ \ \ \ \left(  \text{by the product rule
(\ref{eq.lem.fps.prodrule-fin-fin.eq}), applied to }m_{i}=d\text{ and }%
p_{i,k}=e_{k}\right) \\
&  =\sum_{\left(  k_{1},k_{2},\ldots,k_{n}\right)  \in\left[  d\right]  ^{n}%
}e_{k_{1}}e_{k_{2}}\cdots e_{k_{n}}=\sum_{\left(  x_{1},x_{2},\ldots
,x_{n}\right)  \in\left[  d\right]  ^{n}}e_{x_{1}}e_{x_{2}}\cdots e_{x_{n}}%
\end{align*}
(here, we have renamed the summation index). Summing this over all $\left(
e_{1},e_{2},\ldots,e_{d}\right)  \in\left\{  1,-1\right\}  ^{d}$, we find%
\begin{align*}
&  \sum_{\left(  e_{1},e_{2},\ldots,e_{d}\right)  \in\left\{  1,-1\right\}
^{d}}\left(  e_{1}+e_{2}+\cdots+e_{d}\right)  ^{n}\\
&  =\sum_{\left(  e_{1},e_{2},\ldots,e_{d}\right)  \in\left\{  1,-1\right\}
^{d}}\ \ \sum_{\left(  x_{1},x_{2},\ldots,x_{n}\right)  \in\left[  d\right]
^{n}}e_{x_{1}}e_{x_{2}}\cdots e_{x_{n}}\\
&  =\sum_{\left(  x_{1},x_{2},\ldots,x_{n}\right)  \in\left[  d\right]  ^{n}%
}\ \ \sum_{\left(  e_{1},e_{2},\ldots,e_{d}\right)  \in\left\{  1,-1\right\}
^{d}}e_{x_{1}}e_{x_{2}}\cdots e_{x_{n}}%
\end{align*}
(here, we have interchanged the summation signs). This proves Lemma
\ref{lem.cancel.all-even.l1}.
\end{proof}

Now we shall compute the inner sums on the right hand side of Lemma
\ref{lem.cancel.all-even.l1}:

\begin{lemma}
\label{lem.cancel.all-even.l2}Let $n,d\in\mathbb{N}$. Let $\left(  x_{1}%
,x_{2},\ldots,x_{n}\right)  \in\left[  d\right]  ^{n}$. Then: \medskip

\textbf{(a)} If the $n$-tuple $\left(  x_{1},x_{2},\ldots,x_{n}\right)  $ is
not all-even\footnotemark, then%
\[
\sum_{\left(  e_{1},e_{2},\ldots,e_{d}\right)  \in\left\{  1,-1\right\}  ^{d}%
}e_{x_{1}}e_{x_{2}}\cdots e_{x_{n}}=0.
\]


\textbf{(b)} If the $n$-tuple $\left(  x_{1},x_{2},\ldots,x_{n}\right)  $ is
all-even, then%
\[
\sum_{\left(  e_{1},e_{2},\ldots,e_{d}\right)  \in\left\{  1,-1\right\}  ^{d}%
}e_{x_{1}}e_{x_{2}}\cdots e_{x_{n}}=2^{d}.
\]

\end{lemma}

\footnotetext{See Theorem \ref{thm.cancel.all-even} for the definition of
\textquotedblleft all-even\textquotedblright.}

\begin{proof}
For each $k\in\left[  d\right]  $, let $m_{k}$ be the number of all
$i\in\left[  n\right]  $ satisfying $x_{i}=k$. Then, the $n$-tuple $\left(
x_{1},x_{2},\ldots,x_{n}\right)  $ is all-even if and only if all these
numbers $m_{1},m_{2},\ldots,m_{d}$ are even (by the definition of
\textquotedblleft all-even\textquotedblright).

On the other hand, every $d$-tuple $\left(  e_{1},e_{2},\ldots,e_{d}\right)
\in\mathbb{Z}^{d}$ satisfies%
\begin{equation}
e_{x_{1}}e_{x_{2}}\cdots e_{x_{n}}=e_{1}^{m_{1}}e_{2}^{m_{2}}\cdots
e_{d}^{m_{d}}. \label{pf.lem.cancel.all-even.l2.e=e}%
\end{equation}


\begin{proof}
[Proof of (\ref{pf.lem.cancel.all-even.l2.e=e}):]Let $\left(  e_{1}%
,e_{2},\ldots,e_{d}\right)  \in\mathbb{Z}^{d}$ be any $d$-tuple. Then, each
factor of the product $e_{x_{1}}e_{x_{2}}\cdots e_{x_{n}}$ is one of the $d$
entries $e_{1},e_{2},\ldots,e_{d}$ of this $d$-tuple. Moreover, each given
entry $e_{k}$ appears exactly $m_{k}$ times in the product $e_{x_{1}}e_{x_{2}%
}\cdots e_{x_{n}}$, because $k$ appears exactly $m_{k}$ times among the
subscripts $x_{1},x_{2},\ldots,x_{n}$ (since $m_{k}$ is defined as the number
of all $i\in\left[  n\right]  $ satisfying $x_{i}=k$). Thus, in total, the
product $e_{x_{1}}e_{x_{2}}\cdots e_{x_{n}}$ contains the factor $e_{1}$
exactly $m_{1}$ times, the factor $e_{2}$ exactly $m_{2}$ times, and so on;
therefore, it equals $e_{1}^{m_{1}}e_{2}^{m_{2}}\cdots e_{d}^{m_{d}}$. This
proves (\ref{pf.lem.cancel.all-even.l2.e=e}).
\end{proof}

\textbf{(a)} Assume that the $n$-tuple $\left(  x_{1},x_{2},\ldots
,x_{n}\right)  $ is not all-even. Thus, not all of the numbers $m_{1}%
,m_{2},\ldots,m_{d}$ are even (because the $n$-tuple $\left(  x_{1}%
,x_{2},\ldots,x_{n}\right)  $ is all-even if and only if all these numbers
$m_{1},m_{2},\ldots,m_{d}$ are even). In other words, there exists some
$k\in\left[  d\right]  $ such that $m_{k}$ is odd. Consider this $k$. Hence,
$\left(  -1\right)  ^{m_{k}}=-1$ (since $m_{k}$ is odd).

For each $n$-tuple $\overrightarrow{e}=\left(  e_{1},e_{2},\ldots
,e_{d}\right)  \in\left\{  1,-1\right\}  ^{d}$, we define an integer
\[
\operatorname*{sign}\overrightarrow{e}:=e_{1}^{m_{1}}e_{2}^{m_{2}}\cdots
e_{d}^{m_{d}}.
\]


Let $f:\left\{  1,-1\right\}  ^{d}\rightarrow\left\{  1,-1\right\}  ^{d}$ be
the map that sends each $d$-tuple $\left(  e_{1},e_{2},\ldots,e_{d}\right)  $
to $\left(  e_{1},e_{2},\ldots,e_{k-1},-e_{k},e_{k+1},e_{k+2},\ldots
,e_{d}\right)  $ (in other words, $f$ flips the sign of the $k$-th entry of
the $d$-tuple). This map $f$ is an involution (since flipping the sign of the
$k$-th entry twice results in the restoration of the original sign), and has
no fixed points (since our $d$-tuples $\left(  e_{1},e_{2},\ldots
,e_{d}\right)  $ come from $\left\{  1,-1\right\}  ^{d}$ and thus contain no
zeroes, so that they cannot remain unchanged when we flip the sign of an
entry). Moreover, we have%
\[
\operatorname*{sign}\left(  f\left(  \overrightarrow{e}\right)  \right)
=-\operatorname*{sign}\overrightarrow{e}\ \ \ \ \ \ \ \ \ \ \text{for all
}\overrightarrow{e}\in\left\{  1,-1\right\}  ^{d},
\]
because when we flip the sign of the $k$-th entry $e_{k}$ of a $d$-tuple
$\overrightarrow{e}=\left(  e_{1},e_{2},\ldots,e_{d}\right)  \in\left\{
1,-1\right\}  ^{d}$, the integer $\operatorname*{sign}\overrightarrow{e}%
=e_{1}^{m_{1}}e_{2}^{m_{2}}\cdots e_{d}^{m_{d}}$ is multiplied by $\left(
-1\right)  ^{m_{k}}=-1$. Renaming the variable $\overrightarrow{e}$ as $I$ in
this claim, we can rewrite it as follows:%
\[
\operatorname*{sign}\left(  f\left(  I\right)  \right)  =-\operatorname*{sign}%
I\ \ \ \ \ \ \ \ \ \ \text{for all }I\in\left\{  1,-1\right\}  ^{d}.
\]


Hence, Lemma \ref{lem.sign.cancel2} (applied to $\mathcal{A}=\left\{
1,-1\right\}  ^{d}$ and $\mathcal{X}=\left\{  1,-1\right\}  ^{d}$) yields%
\[
\sum_{I\in\left\{  1,-1\right\}  ^{d}}\operatorname*{sign}I=\sum_{I\in\left\{
1,-1\right\}  ^{d}\setminus\left\{  1,-1\right\}  ^{d}}\operatorname*{sign}%
I=\left(  \text{empty sum}\right)  =0.
\]
But%
\begin{align*}
\sum_{I\in\left\{  1,-1\right\}  ^{d}}\operatorname*{sign}I  &  =\sum_{\left(
e_{1},e_{2},\ldots,e_{d}\right)  \in\left\{  1,-1\right\}  ^{d}}%
\underbrace{\operatorname*{sign}\left(  e_{1},e_{2},\ldots,e_{d}\right)
}_{\substack{=e_{1}^{m_{1}}e_{2}^{m_{2}}\cdots e_{d}^{m_{d}}\\\text{(by the
definition of }\operatorname*{sign}\overrightarrow{e}\text{)}}}\\
&  =\sum_{\left(  e_{1},e_{2},\ldots,e_{d}\right)  \in\left\{  1,-1\right\}
^{d}}\ \ \underbrace{e_{1}^{m_{1}}e_{2}^{m_{2}}\cdots e_{d}^{m_{d}}%
}_{\substack{=e_{x_{1}}e_{x_{2}}\cdots e_{x_{n}}\\\text{(by
(\ref{pf.lem.cancel.all-even.l2.e=e}))}}}\\
&  =\sum_{\left(  e_{1},e_{2},\ldots,e_{d}\right)  \in\left\{  1,-1\right\}
^{d}}e_{x_{1}}e_{x_{2}}\cdots e_{x_{n}}.
\end{align*}
Comparing these two equalities, we find
\[
\sum_{\left(  e_{1},e_{2},\ldots,e_{d}\right)  \in\left\{  1,-1\right\}  ^{d}%
}e_{x_{1}}e_{x_{2}}\cdots e_{x_{n}}=0.
\]
This proves Lemma \ref{lem.cancel.all-even.l2} \textbf{(a)}. \medskip

\textbf{(b)} Assume that the $n$-tuple $\left(  x_{1},x_{2},\ldots
,x_{n}\right)  $ is all-even. Thus, all the numbers $m_{1},m_{2},\ldots,m_{d}$
are even (because the $n$-tuple $\left(  x_{1},x_{2},\ldots,x_{n}\right)  $ is
all-even if and only if all these numbers $m_{1},m_{2},\ldots,m_{d}$ are even).

Now, let $\left(  e_{1},e_{2},\ldots,e_{d}\right)  \in\left\{  1,-1\right\}
^{d}$ be any $d$-tuple consisting of $1$'s and $-1$'s. Then, for each
$k\in\left[  d\right]  $, we have $e_{k}\in\left\{  1,-1\right\}  $ whereas
$m_{k}$ is even (because all the numbers $m_{1},m_{2},\ldots,m_{d}$ are even);
thus we obtain%
\begin{equation}
e_{k}^{m_{k}}=1 \label{pf.lem.cancel.all-even.l2.b.1}%
\end{equation}
(since an even power of $1$ or of $-1$ is always even). Therefore, all the $d$
factors of the product $e_{1}^{m_{1}}e_{2}^{m_{2}}\cdots e_{d}^{m_{d}}$ equal
$1$. Hence, this whole product is a product of $1$'s and thus equals $1$. In
other words, $e_{1}^{m_{1}}e_{2}^{m_{2}}\cdots e_{d}^{m_{d}}=1$. In view of
(\ref{pf.lem.cancel.all-even.l2.e=e}), this rewrites as $e_{x_{1}}e_{x_{2}%
}\cdots e_{x_{n}}=1$.

Forget that we fixed $\left(  e_{1},e_{2},\ldots,e_{d}\right)  $. We thus have
shown that each $\left(  e_{1},e_{2},\ldots,e_{d}\right)  \in\left\{
1,-1\right\}  ^{d}$ satisfies $e_{x_{1}}e_{x_{2}}\cdots e_{x_{n}}=1$.
Therefore,%
\begin{align*}
\sum_{\left(  e_{1},e_{2},\ldots,e_{d}\right)  \in\left\{  1,-1\right\}  ^{d}%
}\underbrace{e_{x_{1}}e_{x_{2}}\cdots e_{x_{n}}}_{=1}  &  =\sum_{\left(
e_{1},e_{2},\ldots,e_{d}\right)  \in\left\{  1,-1\right\}  ^{d}}1=\left\vert
\left\{  1,-1\right\}  ^{d}\right\vert \cdot1\\
&  =\left\vert \left\{  1,-1\right\}  ^{d}\right\vert =\left\vert \left\{
1,-1\right\}  \right\vert ^{d}=2^{d}.
\end{align*}
This proves Lemma \ref{lem.cancel.all-even.l2} \textbf{(b)}.
\end{proof}

We are now ready to prove Theorem \ref{thm.cancel.all-even}:

\begin{proof}
[Proof of Theorem \ref{thm.cancel.all-even}.]Lemma
\ref{lem.cancel.all-even.l1} yields%
\begin{align}
&  \sum_{\left(  e_{1},e_{2},\ldots,e_{d}\right)  \in\left\{  1,-1\right\}
^{d}}\left(  e_{1}+e_{2}+\cdots+e_{d}\right)  ^{n}\nonumber\\
&  =\sum_{\left(  x_{1},x_{2},\ldots,x_{n}\right)  \in\left[  d\right]  ^{n}%
}\ \ \sum_{\left(  e_{1},e_{2},\ldots,e_{d}\right)  \in\left\{  1,-1\right\}
^{d}}e_{x_{1}}e_{x_{2}}\cdots e_{x_{n}}\nonumber\\
&  =\sum_{\substack{\left(  x_{1},x_{2},\ldots,x_{n}\right)  \in\left[
d\right]  ^{n};\\\left(  x_{1},x_{2},\ldots,x_{n}\right)  \text{ is all-even}%
}}\ \ \underbrace{\sum_{\left(  e_{1},e_{2},\ldots,e_{d}\right)  \in\left\{
1,-1\right\}  ^{d}}e_{x_{1}}e_{x_{2}}\cdots e_{x_{n}}}_{\substack{=2^{d}%
\\\text{(by Lemma \ref{lem.cancel.all-even.l2} \textbf{(b)})}}}\nonumber\\
&  \qquad+\sum_{\substack{\left(  x_{1},x_{2},\ldots,x_{n}\right)  \in\left[
d\right]  ^{n};\\\left(  x_{1},x_{2},\ldots,x_{n}\right)  \text{ is not
all-even}}}\ \ \underbrace{\sum_{\left(  e_{1},e_{2},\ldots,e_{d}\right)
\in\left\{  1,-1\right\}  ^{d}}e_{x_{1}}e_{x_{2}}\cdots e_{x_{n}}%
}_{\substack{=0\\\text{(by Lemma \ref{lem.cancel.all-even.l2} \textbf{(a)})}%
}}\nonumber\\
&  =\sum_{\substack{\left(  x_{1},x_{2},\ldots,x_{n}\right)  \in\left[
d\right]  ^{n};\\\left(  x_{1},x_{2},\ldots,x_{n}\right)  \text{ is all-even}%
}}2^{d}+\underbrace{\sum_{\substack{\left(  x_{1},x_{2},\ldots,x_{n}\right)
\in\left[  d\right]  ^{n};\\\left(  x_{1},x_{2},\ldots,x_{n}\right)  \text{ is
not all-even}}}0}_{=0}=\sum_{\substack{\left(  x_{1},x_{2},\ldots
,x_{n}\right)  \in\left[  d\right]  ^{n};\\\left(  x_{1},x_{2},\ldots
,x_{n}\right)  \text{ is all-even}}}2^{d}\nonumber\\
&  =\left(  \text{\# of all all-even }\left(  x_{1},x_{2},\ldots,x_{n}\right)
\in\left[  d\right]  ^{n}\right)  \cdot2^{d}.
\label{sol.count.all-even-tups.c1.pf.b.lhs}%
\end{align}


Now, we shall compute the left hand side of this equality in a different way.
Clearly, a $d$-tuple $\left(  e_{1},e_{2},\ldots,e_{d}\right)  \in\left\{
1,-1\right\}  ^{d}$ is uniquely determined by the set of all its positions
occupied by $-1$'s (that is, by the set $\left\{  i\in\left[  d\right]
\ \mid\ e_{i}=-1\right\}  $). Conversely, for each subset $S$ of $\left[
d\right]  $, there exists a $d$-tuple $\left(  e_{1},e_{2},\ldots
,e_{d}\right)  \in\left\{  1,-1\right\}  ^{d}$ whose set $\left\{  i\in\left[
d\right]  \ \mid\ e_{i}=-1\right\}  $ is $S$. Thus, the map
\begin{align*}
\Phi:\left\{  1,-1\right\}  ^{d}  &  \rightarrow\left\{  \text{subsets of
}\left[  d\right]  \right\}  ,\\
\left(  e_{1},e_{2},\ldots,e_{d}\right)   &  \mapsto\left\{  i\in\left[
d\right]  \ \mid\ e_{i}=-1\right\}
\end{align*}
is a bijection. Moreover, if this bijection $\Phi$ sends a $d$-tuple $\left(
e_{1},e_{2},\ldots,e_{d}\right)  \in\left\{  1,-1\right\}  ^{d}$ to some
subset $S$ of $\left[  d\right]  $, then%
\[
e_{1}+e_{2}+\cdots+e_{d}=d-2\cdot\left\vert S\right\vert
\]
(in fact, each entry $e_{i}$ of the $d$-tuple $\left(  e_{1},e_{2}%
,\ldots,e_{d}\right)  $ equals $-1$ if $i\in S$ and equals $1$ if $i\notin S$;
thus, the sum $e_{1}+e_{2}+\cdots+e_{d}$ has $\left\vert S\right\vert $ many
addends equal to $-1$ while its remaining $d-\left\vert S\right\vert $ addends
are equal to $1$; therefore, this sum totals to $\left\vert S\right\vert
\cdot\left(  -1\right)  +\left(  d-\left\vert S\right\vert \right)
\cdot1=-\left\vert S\right\vert +\left(  d-\left\vert S\right\vert \right)
=d-2\cdot\left\vert S\right\vert $). Thus,%
\[
\sum_{\left(  e_{1},e_{2},\ldots,e_{d}\right)  \in\left\{  1,-1\right\}  ^{d}%
}\left(  e_{1}+e_{2}+\cdots+e_{d}\right)  ^{n}=\sum_{S\subseteq\left[
d\right]  }\left(  d-2\cdot\left\vert S\right\vert \right)  ^{n}%
\]
(since the sum on the left hand side is obtained from the sum on the right
hand side by substituting $\Phi\left(  e_{1},e_{2},\ldots,e_{d}\right)  $ for
$S$). In view of%
\begin{align*}
&  \sum_{S\subseteq\left[  d\right]  }\left(  d-2\cdot\left\vert S\right\vert
\right)  ^{n}\\
&  =\sum_{k=0}^{d}\ \ \sum_{\substack{S\subseteq\left[  d\right]
;\\\left\vert S\right\vert =k}}\left(  d-2\cdot\underbrace{\left\vert
S\right\vert }_{=k}\right)  ^{n}\ \ \ \ \ \ \ \ \ \ \left(
\begin{array}
[c]{c}%
\text{here, we have split the sum}\\
\text{according to the value }\left\vert S\right\vert
\end{array}
\right) \\
&  =\sum_{k=0}^{d}\ \ \underbrace{\sum_{\substack{S\subseteq\left[  d\right]
;\\\left\vert S\right\vert =k}}\left(  d-2k\right)  ^{n}}_{=\left(  \text{\#
of all }S\subseteq\left[  d\right]  \text{ satisfying }\left\vert S\right\vert
=k\right)  \cdot\left(  d-2k\right)  ^{n}}\\
&  =\sum_{k=0}^{d}\underbrace{\left(  \text{\# of all }S\subseteq\left[
d\right]  \text{ satisfying }\left\vert S\right\vert =k\right)  }_{=\left(
\text{\# of all }k\text{-element subsets of }\left[  d\right]  \right)
=\dbinom{d}{k}}\cdot\left(  d-2k\right)  ^{n}\\
&  =\sum_{k=0}^{d}\dbinom{d}{k}\left(  d-2k\right)  ^{n},
\end{align*}
we can rewrite this as%
\[
\sum_{\left(  e_{1},e_{2},\ldots,e_{d}\right)  \in\left\{  1,-1\right\}  ^{d}%
}\left(  e_{1}+e_{2}+\cdots+e_{d}\right)  ^{n}=\sum_{k=0}^{d}\dbinom{d}%
{k}\left(  d-2k\right)  ^{n}.
\]
Comparing this with \eqref{sol.count.all-even-tups.c1.pf.b.lhs}, we obtain%
\[
\left(  \text{\# of all all-even }\left(  x_{1},x_{2},\ldots,x_{n}\right)
\in\left[  d\right]  ^{n}\right)  \cdot2^{d}=\sum_{k=0}^{d}\dbinom{d}%
{k}\left(  d-2k\right)  ^{n}.
\]
Dividing this equality by $2^{d}$, we obtain%
\[
\left(  \text{\# of all all-even }\left(  x_{1},x_{2},\ldots,x_{n}\right)
\in\left[  d\right]  ^{n}\right)  =\dfrac{1}{2^{d}}\sum_{k=0}^{d}\dbinom{d}%
{k}\left(  d-2k\right)  ^{n}.
\]
This proves Theorem \ref{thm.cancel.all-even}.
\end{proof}

\begin{noncompile}
Compare with \cite[Corollary 2.5]{Stanle18}.
\end{noncompile}
